% Part: propositional-logic
% Chapter: syntax-and-semantics
% Section: valuations-sat

\documentclass[../../../include/open-logic-section]{subfiles}

\begin{document}

\olfileid{pl}{syn}{val}

\olsection{\usetoken{P}{valuation} lan Panyukupan}

\begin{defn}[!!^{valuation}s]
$\{\True, \False\}$ dadi himpunan rong nilai kabeneran, ``bener'' lan
``salah''. \emph{!!{valuation}} kanggo $\Lang{L_0}$ yaiku
fungsi~$\pAssign{v}$ kang masangake $\True$ utawa $\False$ marang
!!{propositional variable}s ing basa kasebut, yaiku $\pAssign{v} \colon
\PVar \to \{\True, \False \}$.
\end{defn}

\begin{defn}
  Yen diwenehi !!a{valuation}~$\pAssign{v}$, temtokna fungsi evaluasi
  $\pValue{v} \colon \Frm[L_0] \to \{\True, \False \}$ kanthi induktif:
  \begin{align*}
    \iftag{prvFalse}{\pValue{v}(\lfalse) & = \False; \\}{}
    \iftag{prvTrue}{\pValue{v}(\ltrue) & = \True; \\}{}
    \pValue{v}(\Obj p_n) & = \pAssign{v}(\Obj p_n); \\
    \iftag{prvNot}{\pValue{v}(\lnot !A) & = \begin{cases}
        \True & \text{yen } \pValue{v}(!A) = \False;\\
        \False & \text{yen ora.}
      \end{cases} \\ }{}
    \iftag{prvAnd}{\pValue{v}(!A \land !B) & = \begin{cases}
        \True &
        \text{yen $\pValue{v}(!A) = \True$ lan $\pValue{v}(!B) = \True$;}\\
        \False &
        \text{yen $\pValue{v}(!A) = \False$ utawa $\pValue{v}(!B) = \False$}.
      \end{cases}\\}{}
    \iftag{prvOr}{\pValue{v}(!A \lor !B) & = \begin{cases}
        \True &
        \text{yen $\pValue{v}(!A) = \True$ utawa $\pValue{v}(!B) = \True$;}\\
        \False &
        \text{yen $\pValue{v}(!A) = \False$ lan $\pValue{v}(!B) = \False$}.
      \end{cases}\\}{}
    \iftag{prvIf}{\pValue{v}(!A \lif !B) & = \begin{cases}
        \True &
        \text{yen $\pValue{v}(!A) = \False$ utawa $\pValue{v}(!B) = \True$;}\\
        \False &
        \text{yen $\pValue{v}(!A) = \True$ lan $\pValue{v}(!B) = \False$}.
      \end{cases}\\}{}
    \iftag{prvIff}{\pValue{v}(!A \liff !B) & = \begin{cases}
        \True &
        \text{yen $\pValue{v}(!A) = \pValue{v}(!B)$;}\\
        \False &
        \text{yen $\pValue{v}(!A) \neq \pValue{v}(!B)$}.
      \end{cases}}{}
\end{align*}
\end{defn}

\begin{explain}
Klausa-klausa kasebut cocog karo tabel kabeneran iki:
\begin{center}
\iftag{prvNot}{
  \begin{tabular}{|c||c|} \hline
    $!A$ & $ \lnot !A$ \\
    \hline \hline
    $\True$ & $\False$ \\
    $\False$ & $\True$ \\
    \hline
  \end{tabular}
}{}
\iftag{prvAnd}{
  \begin{tabular}{|cc||c|} \hline
    $!A$ & $!B$ & $!A \land !B$ \\
    \hline \hline
    $\True$ & $\True$ & $\True$ \\
    $\True$ & $\False$ & $\False$ \\
    $\False$ & $\True$ & $\False$ \\
    $\False$ & $\False$ & $\False$ \\
    \hline
  \end{tabular}
}{}
\iftag{prvOr}{
  \begin{tabular}{|cc||c|} \hline
    $!A$ & $!B$ & $!A \lor !B$ \\
    \hline \hline
    $\True$ & $\True$ & $\True$ \\
    $\True$ & $\False$ & $\True$ \\
    $\False$ & $\True$ & $\True$ \\
    $\False$ & $\False$ & $\False$ \\
    \hline
  \end{tabular}
}{}%
\iftag{notprvNot,notprvAnd,notprvOr,notprvIf}{}{\\[1em]}
\iftag{prvIf}{
  \begin{tabular}{|cc||c|} \hline
    $!A$ & $!B$ & $!A \lif !B$ \\
    \hline \hline
    $\True$ & $\True$ & $\True$ \\
    $\True$ & $\False$ & $\False$ \\
    $\False$ & $\True$ & $\True$ \\
    $\False$ & $\False$ & $\True$ \\
    \hline
  \end{tabular}
}{}
\iftag{prvIff}{
  \begin{tabular}{|cc||c|} \hline
    $!A$ & $!B$ & $!A \liff !B$ \\
    \hline \hline
    $\True$ & $\True$ & $\True$ \\
    $\True$ & $\False$ & $\False$ \\
    $\False$ & $\True$ & $\False$ \\
    $\False$ & $\False$ & $\True$ \\
    \hline
  \end{tabular}
}{}
\end{center}
\end{explain}

\begin{prob}
Gatekna yen ing $\Lang{L_0}$ ditambahake konektif ternèr
$\diamondsuit$ kanthi evaluasi kang diwenehake dening
\begin{gather*}
  \pValue{v}(\diamondsuit( !A , !B , !C ) ) = \begin{cases}
    \pValue{v}( !B ) &
    \text{yen $\pValue{v}(!A) = \True$;}\\
    \pValue{v}( !C ) &
    \text{yen $\pValue{v}(!A) = \False $}.
  \end{cases}
\end{gather*}
Tulisen tabel kabeneran kanggo konektif iki.
\end{prob}

\begin{thm}[Penentuan Lokal]
\ollabel{thm:LocalDetermination} Upamakna $\pAssign{v_1}$ lan
$\pAssign{v_2}$ iku !!{valuation}s kang menehi nilai padha marang
!!{propositional
variable}s kang muncul ing $!A$, yaiku
$\pAssign{v_1}(\Obj p_n) = \pAssign{v_2}(\Obj p_n)$ saben $\Obj p_n$
muncul ing sawijining !!{formula}~$!A$. Banjur $\pValue{v_1}$ lan
$\pValue{v_2}$ uga menehi nilai padha marang~$!A$, yaiku
$\pValue{v_1}(!A) = \pValue{v_2}(!A)$.
\end{thm}

\begin{proof}
Kanthi induksi ing $!A$.
\end{proof}

\begin{defn}[Panyukupan]
\ollabel{defn:satisfaction} Kita bisa nemtokake kanthi induktif gagasan
  \emph{panyukupan !!a{formula}~$!A$ dening
  !!a{valuation}~$\pAssign{v}$}, $\pSat{v}{!A}$, kaya mangkene.
  (Kita nulis $\pSat/{v}{!A}$ kanthi teges ``ora $\pSat{v}{!A}$.'' )
\begin{enumerate}
\tagitem{prvFalse}{%
  \indcase{!A}{\lfalse}{$\pSat/{v}{\indfrm}$.}}{}

\tagitem{prvTrue}{%
  \indcase{!A}{\ltrue}{$\pSat{v}{\indfrm}$.}}{}

\item \indcase{!A}{\Obj p_i}{$\pSat{v}{\indfrm}$
  yen lan mung yen $\pAssign{v}(\Obj p_i) = \True$.}

\tagitem{prvNot}{%
  \indcase{!A}{\lnot !B}{$\pSat{v}{\indfrm}$ yen lan mung yen
    $\pSat/{v}{!B}$.}}{}

\tagitem{prvAnd}{%
  \indcase{!A}{(!B \land !C)}{$\pSat{v}{\indfrm}$ yen lan mung yen
    $\pSat{v}{!B}$ lan $\pSat{v}{!C}$.}}{}

\tagitem{prvOr}{%
  \indcase{!A}{(!B \lor !C)}{$\pSat{v}{\indfrm}$ yen lan mung yen
    $\pSat{v}{!B}$ utawa $\pSat{v}{!C}$ (utawa kalorone).}}{}

\tagitem{prvIf}{%
  \indcase{!A}{(!B \lif !C)}{$\pSat{v}{\indfrm}$ yen lan mung yen
    $\pSat/{v}{!B}$ utawa $\pSat{v}{!C}$ (utawa kalorone).}}{}

\tagitem{prvIff}{%
  \indcase{!A}{(!B \liff !C)}{$\pSat{v}{\indfrm}$ yen lan mung yen
    $\pSat{v}{!B}$ lan $\pSat{v}{!C}$ kalorone, utawa $\pSat{v}{!B}$
    lan $\pSat{v}{!C}$ kalorone ora lumaku.}}{}
\end{enumerate}
Yen $\Gamma$ iku himpunan !!{formula}s, $\pSat{v}{\Gamma}$ yen lan mung
yen $\pSat{v}{!A}$ kanggo saben~$!A \in \Gamma$.
\end{defn}

\begin{prop}\ollabel{prop:sat-value}
  $\pSat{v}{!A}$ yen lan mung yen $\pValue{v}(!A) = \True$.
\end{prop}

\begin{proof}
  Kanthi induksi ing~$!A$.
\end{proof}

\begin{prob}
  Buktekna \olref[pl][syn][val]{prop:sat-value}
\end{prob}

\end{document}
